\section{Statements and assumptions}\label{sec:statements}

For $n\geq1$, let $\Omega(n)$ count prime factors with multiplicity.
For each prime $p\geq3$, list the prime factors of $p-1$ in decreasing
order, with multiplicity, as $q_1(p),q_2(p),\ldots$, and put
\[
 V_i(p)=\frac{\log q_i(p)}{\log(p-1)}.
\]
After the last prime factor, set $V_i(p)=0$. The sequence $V(p)$ is a
decreasing partition of one. Let $U_1,U_2,\ldots$ be independent uniform
variables on $(0,1)$, and define
\[
 T_i=U_i\prod_{j<i}(1-U_j).
\]
The sequence $(T_i)$ has the $\GEM(1)$ distribution. Its decreasing
rearrangement $L=(L_1,L_2,\ldots)$ has the Poisson--Dirichlet distribution
$\PD(1)$. We write $\pi(X;M,b)$ for the number of primes at most $X$
congruent to $b$ modulo $M$.

The distribution of $V(p)$ alone does not determine the parity of
$\Omega(p-1)$, or its residue modulo a larger integer. Small prime factors
can change that residue without producing a detectable change in finitely
many normalized factor sizes. We therefore prove cancellation of each
nontrivial root-of-unity weight against marked factor statistics before
passing to a probability limit.

\begin{assumption}[Analytic source inputs]\label{ass:inputs}
Assume the specified results and proof estimates in the following
manuscripts, at revision
\texttt{fd4aeeb2ee4fc729c18d98444fed42fd0529eeeb} of the OpenAI
mathematics repository:
\begin{enumerate}[label=\textup{(\roman*)}]
\item the graph transference and residual ideal-family inputs, comparison
estimates, long-prime estimates, and positive extraction estimates in
\emph{The Poisson--Dirichlet law for prime predecessors} \cite{D};
\item the general physical pairing and ideal-kernel results in
\emph{Weighted dilation graphs, smooth shifted primes, and totient fibers}
\cite{S}, as restated in \cite[Section 6]{D};
\item the scalar coefficient criterion in
\emph{Primitive roots for every admissible integer base}
\cite[Section 8]{P}, including the scalar Cauchy-square estimate in its proof;
\item the zero-free half-plane $\Rea s>7/8$ for Dirichlet $L$-functions
in \emph{The quasi-Riemann hypothesis} \cite{Z}, with the principal pole
at $s=1$ retained.
\end{enumerate}
The precise source labels, hypotheses needed here, and uses are recorded
in Section~\ref{sec:sources}. These source results are premises of this paper.
\end{assumption}

The fixed-progression Poisson--Dirichlet law is proved below as
Proposition~\ref{prop:unsigned}; it is not inferred by conditioning the
unrestricted source law. Its progression argument is an explicit
dependency of both the residue theorem and the squarefree refinement.

\begin{theorem}[Factor-count residues]\label{thm:main}
Assume Assumption~\ref{ass:inputs}. Fix integers $M,m,k\geq1$, a reduced
class $b\pmod M$, a residue $a\pmod m$, and a continuous function
$H:[0,1]^k\to\C$. Then
\begin{equation}\label{eq:main}
 \frac{1}{\pi(X;M,b)}
 \sum_{\substack{3\leq p\leq X\\p\equiv b\ (M)}}
 \ind_{\Omega(p-1)\equiv a\ (m)}H(V_1(p),\ldots,V_k(p))
 \longrightarrow \frac1m\E H(L_1,\ldots,L_k).
\end{equation}
The same limit holds for primes in $(x,2x]$, with denominator
$\pi(2x;M,b)-\pi(x;M,b)$.
\end{theorem}

Every displayed parameter is fixed before the real endpoint tends to
infinity. The theorem gives no rate and no uniformity in growing moduli.
Taking $H=1$ gives uniform residues; dividing by their limiting
probability $1/m$ gives the conditional factor-size law.

Write $\lambda(n)=(-1)^{\Omega(n)}$ and let $\mu$ be the M\"obius
function. For a prime $q$, set $e_q=v_q(M)$ and define
\begin{equation}\label{eq:rho}
 \rho_q(M,b)=
 \begin{cases}
  1/[q(q-1)],&e_q=0,\\
  1/q,&e_q=1,\ b\equiv1\pmod q,\\
  0,&e_q=1,\ b\not\equiv1\pmod q,\\
  1,&e_q\geq2,\ b\equiv1\pmod{q^2},\\
  0,&e_q\geq2,\ b\not\equiv1\pmod{q^2}.
 \end{cases}
 \qquad C_{M,b}=\prod_q(1-\rho_q(M,b)).
\end{equation}
The product includes $q=2$.

\begin{theorem}[Parity and squarefreeness]\label{thm:squarefree}
Under Assumption~\ref{ass:inputs}, if $p\geq3$ is uniform among the class
primes in either sampling convention of Theorem~\ref{thm:main}, then
\[
 (V_1(p),\ldots,V_k(p),\lambda(p-1),\mu^2(p-1))
 \Longrightarrow\ (L_1,\ldots,L_k,S,B),
\]
where $S$ is a fair sign, $B$ is Bernoulli with parameter $C_{M,b}$,
and $L,S,B$ are mutually independent. In particular, for $s\in\{-1,1\}$,
\begin{align}
 \E[\mu(p-1)H(V_1(p),\ldots,V_k(p))]&\longrightarrow0,\label{eq:mu}\\
 \E[\ind_{\mu(p-1)=s}H(V_1(p),\ldots,V_k(p))]
   &\longrightarrow\frac{C_{M,b}}2\E H(L_1,\ldots,L_k).\notag
\end{align}
\end{theorem}

The notation $\E$ on the left denotes the finite prime average.
If a square is forced by the class, $C_{M,b}=0$ and the Bernoulli
variable is zero. Conditioning on squarefreeness or on a M\"obius sign
requires $C_{M,b}>0$. For $M=1$, each M\"obius sign has density
$\tfrac12\prod_q(1-1/[q(q-1)])$ relative to the primes.

For a finite set $\mathcal S$ of primes, define
\[
 \mu_{\mathcal S}(n)=\mu\left(n\Big/\prod_{q\in\mathcal S}q^{v_q(n)}\right).
\]
Thus square factors in $\mathcal S$ do not force this function to vanish.

\begin{corollary}[Prescribed valuations]\label{cor:moree}
Under Assumption~\ref{ass:inputs}, for every finite $\mathcal S$ and
fixed nonnegative integers $e_q$, $q\in\mathcal S$,
\[
 \sum_{\substack{p\leq X\\v_q(p-1)=e_q\ (q\in\mathcal S)}}
      \mu_{\mathcal S}(p-1)=o(\pi(X)).
\]
This is the parity assertion of Moree and Hommersom
\cite[Conjecture 1]{MH}, under the stated source assumptions.
\end{corollary}

\subsection{Related work and the proof's scope}

The ordinary-integer distribution results of Pintz \cite{Pintz} and
Granville and Shao \cite{GS} supply progression estimates used below.
They do not restrict the averaging variable to prime predecessors.
Lichtman \cite{Lichtman} proves cancellation averaged over a growing
range of shifts. Carella \cite[Theorems 1.1--1.2]{Carella} states
shifted-prime cancellation subject to a separate distribution hypothesis.
These statements must be distinguished from the fixed-shift joint law
in \eqref{eq:main}. We make no claim to an exhaustive priority search.

The factor-count residue is recovered by finite Fourier inversion.
The analytic problem is to prove cancellation for
$f_z(n)=z^{\Omega(n)}$, for each nontrivial root of unity $z$.
For $z=-1$, the ordinary mean has arbitrary logarithmic cancellation.
For nonreal $z$, the Selberg--Delange expansion has a nonzero principal
term. Sections~\ref{sec:model} and~\ref{sec:extraction} construct and
propagate that term. Section~\ref{sec:probability} derives the joint
laws from the resulting interior factor estimates.
